\documentclass[aspectratio=169,10pt]{beamer}

\usetheme{Madrid}
\usecolortheme{whale}

\usepackage[utf8]{inputenc}
\usepackage[spanish,es-tabla]{babel}
\usepackage{amsmath,amssymb}
\usepackage{booktabs}
\usepackage{graphicx}

% Configuración de colores institucionales UniNorte
\definecolor{UniNorteBlue}{RGB}{20, 50, 95}
\definecolor{UniNorteRed}{RGB}{180, 40, 40}
\definecolor{UniNorteGreen}{RGB}{16, 120, 60}

\setbeamercolor{palette primary}{bg=UniNorteBlue,fg=white}
\setbeamercolor{palette secondary}{bg=UniNorteBlue!85,fg=white}
\setbeamercolor{palette tertiary}{bg=UniNorteBlue!70,fg=white}
\setbeamercolor{titlelike}{parent=palette primary}
\setbeamercolor{block title}{bg=UniNorteBlue,fg=white}
\setbeamercolor{block body}{bg=UniNorteBlue!10}

\title[\textbf{Grafos y Perceptrón en Cra 51B}]{\textbf{Optimización del Despacho de Refuerzo en Redes Urbanas Mediante Teoría de Grafos y Perceptrón de Rosenblatt}}
\subtitle{Caso de Estudio: Corredor Universitario de la Cra 51B (Obras de ``La Gran Vía'')\\
\small \textbf{Arquitectura de Doble Escala (6 Nodos Conocidos + 15 Nodos Reales) y Dashboard Interactivo}}
\author[\textbf{MAT 4021}]{\textbf{Equipo de Investigación Estudiantil}\\Estudiante 1 \and Estudiante 2 \and Estudiante 3}
\institute[UniNorte]{Departamento de Matemáticas y Estadística\\Universidad del Norte\\MAT 4021 - Matemáticas Discretas (MSc. Iván Gómez)}
\date{Noviembre 2026}

\begin{document}

% -------------------------------------------------------------
% SLIDE 1: Portada
% -------------------------------------------------------------
\begin{frame}
    \titlepage
\end{frame}

% -------------------------------------------------------------
% SLIDE 2: Problema y Justificación Territorial (Estudiante 1)
% -------------------------------------------------------------
\begin{frame}{Contexto Territorial: Impacto Vial de ``La Gran Vía''}
    \begin{columns}[T]
        \begin{column}{0.55\textwidth}
            \textbf{Problemática en la Carrera 51B:}
            \begin{itemize}
                \item Movilidad diaria: \textbf{$>45.000$ usuarios} (UniNorte, UniAtlántico, UniLibre, colegios).
                \item Reducción severa de calzadas por obras civiles de modernización vial.
                \item Tiempos de espera en paradero: escalaron de $12$--$15$ min a \textbf{$35$--$50$ min} en horas pico.
                \item Sobrecupo crítico en flotas: Sobusa, Coolitoral y Transmetro U30 (tasa de sobrecupo $>65\%$).
            \end{itemize}
        \end{column}
        \begin{column}{0.45\textwidth}
            \begin{block}{Objetivo de Decisión Discreta}
                Determinar analíticamente qué paraderos demandan \textbf{despacho urgente de buses de refuerzo en vacío ($y=1$)} frente a régimen estándar ($y=0$), sin recurrir a cajas negras de IA.
            \end{block}
            \vspace{0.2cm}
            \begin{exampleblock}{Metodología de Doble Escala}
                \begin{itemize}
                    \item \textbf{Subgrafo 6 Nodos:} Prueba con resultado conocido (Rúbrica).
                    \item \textbf{Red 15 Nodos:} Modelo real de todo el corredor y simulación de cierres viales.
                \end{itemize}
            \end{exampleblock}
        \end{column}
    \end{columns}
\end{frame}

% -------------------------------------------------------------
% SLIDE 3: Red Escalada del Corredor (15 Nodos) y Cierres (Estudiante 1)
% -------------------------------------------------------------
\begin{frame}{Red Arterial Escalada ($n=15$) y Simulación de Cierre Vial}
    \begin{columns}[T]
        \begin{column}{0.52\textwidth}
            \textbf{Modelo Arterial Completo (Cra 51B):}
            \begin{itemize}
                \item 15 macro-paraderos reales: Circunvalar ($P01$) a Pto. Colombia ($P15$).
                \item Incorpora: Makro, Portoazul, San José, UniNorte (P2 y P4), Británico, UniLibre, UAN, UniAtlántico, Campestre.
                \item 20 aristas base que modelan calzadas principales y desvíos vecinales.
            \end{itemize}
            \begin{alertblock}{Simulación de Cierre en La Gran Vía}
                Al clausurar la arista $(P04, P07)$ (Portoazul $\leftrightarrow$ UniNorte P4):
                \begin{itemize}
                    \item El flujo se desvía por Colegio San José ($P05$).
                    \item Su intermediación salta de $C_B'=0.016$ a $\mathbf{0.240}$ ($+0.224$).
                \end{itemize}
            \end{alertblock}
        \end{column}
        \begin{column}{0.48\textwidth}
            \textbf{Comportamiento Topológico ante Cierres:}
            \centering
            \small
            \begin{tabular}{lccc}
                \toprule
                \textbf{Paradero} & \textbf{$C_B'$ Normal} & \textbf{$C_B'$ Cierre} & \textbf{Efecto} \\
                \midrule
                P02 (Makro) & 0.159 & 0.165 & Arterial \\
                P03 (Le Champ) & 0.022 & 0.090 & Desvío $\uparrow$ \\
                P04 (Portoazul) & 0.225 & 0.163 & Bloqueo $\downarrow$ \\
                P05 (San José) & 0.016 & \textbf{0.240} & \textbf{Saturación} $\uparrow$ \\
                P06 (UniNorte P2) & 0.104 & 0.167 & Bypass $\uparrow$ \\
                P07 (UniNorte P4) & 0.489 & 0.377 & Demanda \\
                \bottomrule
            \end{tabular}
            \vspace{0.2cm}
            \footnotesize \textit{Demuestra algebraicamente la congestión en rutas de desvío.}
        \end{column}
    \end{columns}
\end{frame}

% -------------------------------------------------------------
% SLIDE 4: Caso de Prueba Conocido - 6 Nodos (Estudiante 1)
% -------------------------------------------------------------
\begin{frame}{Caso de Prueba Conocido: Subgrafo Diamante ($|V|=6, |E|=6$)}
    \begin{columns}[T]
        \begin{column}{0.50\textwidth}
            \textbf{Subgrafo Analítico de Verificación:}
            \begin{itemize}
                \item $v_1$: Club Campestre (Acceso Norte)
                \item $v_2$: Clínica Portoazul (Tramo Gran Vía)
                \item $v_3$: UniNorte Puerta 4 (Demanda Máx.)
                \item $v_4$: Makro (Bypass / Articulación)
                \item $v_5$: UniAtlántico (Nodo Secundario)
                \item $v_6$: Retorno Sur (Terminal Enlace)
            \end{itemize}
            \textit{Estructura Diamante:} $\{v_1, v_2, v_3, v_4\}$ modela el desvío paralelo por obras.
        \end{column}
        \begin{column}{0.50\textwidth}
            \textbf{Matriz de Adyacencia Binaria $A(G)$:}
            $$A = \begin{pmatrix}
            0 & 1 & 1 & 0 & 0 & 0 \\
            1 & 0 & 0 & 1 & 0 & 0 \\
            1 & 0 & 0 & 1 & 0 & 0 \\
            0 & 1 & 1 & 0 & 1 & 0 \\
            0 & 0 & 0 & 1 & 0 & 1 \\
            0 & 0 & 0 & 0 & 1 & 0
            \end{pmatrix}$$
            \small Simétrica con diagonal nula; grado de $v_4 = 3$ (nodo articulador).
        \end{column}
    \end{columns}
\end{frame}

% -------------------------------------------------------------
% SLIDE 5: Algoritmo BFS por Capas y Distancias (Estudiante 2)
% -------------------------------------------------------------
\begin{frame}{Algoritmo BFS por Capas y Matriz de Distancias $D$}
    \begin{columns}[T]
        \begin{column}{0.50\textwidth}
            \textbf{Búsqueda en Anchura (BFS):}
            \begin{itemize}
                \item Partición por capas desde cada raíz $s$:
                $$L_k = \{u \in V \mid d(s, u) = k\}$$
                \item Al ser un grafo no ponderado, BFS garantiza el cálculo determinista del camino geodésico exacto.
                \item Preservación de predecesores múltiples para rastrear concurrencia de caminos.
            \end{itemize}
        \end{column}
        \begin{column}{0.50\textwidth}
            \textbf{Matriz de Distancias Geodésicas $D$:}
            $$D = \begin{pmatrix}
            0 & 1 & 1 & 2 & 3 & 4 \\
            1 & 0 & 2 & 1 & 2 & 3 \\
            1 & 2 & 0 & 1 & 2 & 3 \\
            2 & 1 & 1 & 0 & 1 & 2 \\
            3 & 2 & 2 & 1 & 0 & 1 \\
            4 & 3 & 3 & 2 & 1 & 0
            \end{pmatrix}$$
        \end{column}
    \end{columns}
\end{frame}

% -------------------------------------------------------------
% SLIDE 6: Intermediación de Freeman y Empates (Estudiante 2)
% -------------------------------------------------------------
\begin{frame}{Centralidad de Intermediación ($C_B$) y Factor de Normalización}
    \textbf{Fórmula de Freeman para Grafos No Dirigidos:}
    $$C_B(v) = \sum_{s < t,\; s \neq v \neq t} \frac{\sigma_{st}(v)}{\sigma_{st}}$$

    \begin{columns}[T]
        \begin{column}{0.52\textwidth}
            \begin{block}{Caminos Empatados con Peso Fraccionario}
                Entre $v_1$ (Campestre) y $v_4$ (Makro) existen dos geodésicas de longitud 2:
                \begin{itemize}
                    \item Ruta A: $v_1 \to v_2 \to v_4$ \quad ($\sigma_{v_1 v_4} = 2$)
                    \item Ruta B: $v_1 \to v_3 \to v_4$
                \end{itemize}
                Aporte para $v_2$ y $v_3$: $\frac{\sigma(v)}{\sigma} = \frac{1}{2} = \mathbf{0.50}$.
                \vspace{0.1cm}
                \footnotesize \textit{Total: 8 fracciones de 0.50 en la red (pares $v_1$-$v_4$, $v_1$-$v_5$, $v_1$-$v_6$, $v_2$-$v_3$) que cierran la suma exacta a $C_B$.}
            \end{block}
        \end{column}
        \begin{column}{0.48\textwidth}
            \begin{block}{Normalización Combinatoria ($C_B'$)}
                Pares posibles sin incluir el nodo evaluado:
                $$\binom{n - 1}{2} = \binom{5}{2} = \frac{5 \times 4}{2} = \mathbf{10}$$
                $$C_B'(v) = \frac{C_B(v)}{10} \in [0, 1]$$
                Makro ($v_4$) obtiene $C_B'=0.65$ (máximo puente).
            \end{block}
        \end{column}
    \end{columns}
\end{frame}

% -------------------------------------------------------------
% SLIDE 7: Tabla Consolidada de Métricas (Estudiante 2)
% -------------------------------------------------------------
\begin{frame}{Tabla Consolidada de Centralidades (Caso 6 Nodos)}
    \centering
    \small
    \begin{tabular}{clcccccc}
        \toprule
        \textbf{Vértice} & \textbf{Paradero} & \textbf{$C_D$} & \textbf{$\sum d$} & \textbf{$C_C$} & \textbf{$C_B$} & \textbf{$C_B'=C_B/10$} & \textbf{Afluencia ($x_2$)} \\
        \midrule
        $v_1$ & Club Campestre & 2 & 11 & 0.455 & 0.50 & \textbf{0.05} & 0.20 \\
        $v_2$ & Clínica Portoazul & 2 & 9 & 0.556 & 1.50 & \textbf{0.15} & 0.70 \\
        $v_3$ & UniNorte Puerta 4 & 2 & 9 & 0.556 & 1.50 & \textbf{0.15} & 0.95 \\
        $v_4$ & Makro / Transversal & 3 & 7 & 0.714 & 6.50 & \textbf{0.65} & 0.80 \\
        $v_5$ & UniAtlántico / Colegios & 2 & 9 & 0.556 & 4.00 & \textbf{0.40} & 0.55 \\
        $v_6$ & Retorno Sur 51B & 1 & 13 & 0.385 & 0.00 & \textbf{0.00} & 0.15 \\
        \bottomrule
    \end{tabular}
    \vspace{0.4cm}
    \begin{alertblock}{Reto Analítico en el Espacio de Características $\mathbb{R}^2$}
        Vector de entrada: $\mathbf{x} = [C_B', \text{Afluencia}_{norm}]^T$. ¿Cómo clasificar de forma determinista y sin cajas negras qué paraderos requieren refuerzo urgente?
    \end{alertblock}
\end{frame}

% -------------------------------------------------------------
% SLIDE 8: Perceptrón de Rosenblatt (Estudiante 3)
% -------------------------------------------------------------
\begin{frame}{Modelo de IA Discreta: Perceptrón de Rosenblatt (1958)}
    \begin{columns}[T]
        \begin{column}{0.50\textwidth}
            \textbf{Formulación Matemática:}
            \begin{itemize}
                \item Vector de entrada: $\mathbf{x} = [C_B', \text{Afluencia}_{norm}]^T$.
                \item Combinación lineal: $z = w_1 x_1 + w_2 x_2 + b$.
                \item Función Heaviside: $\hat{y} = \phi(z) = \begin{cases} 1, & z \ge 0 \\ 0, & z < 0 \end{cases}$
            \end{itemize}
            \textbf{Condición Inicial \& Tasa:}
            \begin{itemize}
                \item $\mathbf{w}^{(0)} = [0, 0]^T, \; b^{(0)} = -1.0$ (inercia operativa $\theta=1.0$).
                \item Tasa de aprendizaje: $\eta = 0.5$.
            \end{itemize}
        \end{column}
        \begin{column}{0.50\textwidth}
            \textbf{Regla de Actualización Discreta:}
            $$e = y - \hat{y} \in \{-1, 0, 1\}$$
            $$\Delta \mathbf{w} = \eta \cdot e \cdot \mathbf{x}, \quad \Delta b = \eta \cdot e$$
            $$\mathbf{w}^{(k+1)} = \mathbf{w}^{(k)} + \Delta \mathbf{w}, \quad b^{(k+1)} = b^{(k)} + \Delta b$$
            \vspace{0.2cm}
            \textit{Cero cajas negras: álgebra de matrices y umbrales discretos puros en NumPy.}
        \end{column}
    \end{columns}
\end{frame}

% -------------------------------------------------------------
% SLIDE 9: Prueba de Escritorio y Convergencia (Estudiante 3)
% -------------------------------------------------------------
\begin{frame}{Prueba de Escritorio: Convergencia Exacta en Época 2}
    \begin{columns}[T]
        \begin{column}{0.55\textwidth}
            \textbf{Dinámica de Aprendizaje Demostrada:}
            \begin{itemize}
                \item \textbf{Época 1:} 3 errores de clasificación ($p_3, p_4, p_5$).
                \item Pesos al final de Época 1:
                $$\mathbf{w} = [0.20, 0.60]^T, \quad b = -0.50$$
                \item \textbf{Época 2:} Se evalúan los 6 paraderos.
                \item \textbf{Error Total en Época 2 = 0 (Convergencia alcanzada)}.
            \end{itemize}
        \end{column}
        \begin{column}{0.45\textwidth}
            \begin{block}{Hiperplano Separador Aprendido}
                $$0.20 C_B' + 0.60 \text{Afluencia} - 0.50 = 0$$
                $$\text{Afluencia} = -0.3333 \cdot C_B' + 0.8333$$
                \begin{itemize}
                    \item $w_1 = 0.20 > 0$: Topología aporta al refuerzo.
                    \item $w_2 = 0.60 > 0$: Demanda aporta al refuerzo.
                    \item Pendiente negativa: Makro ($C_B'=0.65$) requiere menor afluencia para disparar refuerzo.
                \end{itemize}
            \end{block}
        \end{column}
    \end{columns}
\end{frame}

% -------------------------------------------------------------
% SLIDE 10: Demostración Interactiva y Métricas de Impacto (Equipo)
% -------------------------------------------------------------
\begin{frame}{Demostración en Vivo: Aplicación Web y Métricas de Impacto}
    \begin{columns}[T]
        \begin{column}{0.50\textwidth}
            \textbf{Consola Web Interactiva (\texttt{localhost:8080}):}
            \begin{itemize}
                \item \textbf{Pestaña 1:} Mapa interactivo Leaflet (OSM y Esri Satelital) de 15 paraderos con toggle en vivo de cierre vial en La Gran Vía.
                \item \textbf{Pestaña 2:} Caso conocido (6 nodos) con SVG dinámico por capas BFS, matrices $A, D$ y tabla auditable de los 8 empates ($0.50$).
                \item \textbf{Pestaña 3:} Sliders de $\eta, b_0$, espacio $\mathbb{R}^2$ con recta de separación y traza aritmética por época.
                \item \textbf{Pestaña 4:} Simulador configurable con reducción del 58.8\% en espera y horas-pasajero ahorradas.
            \end{itemize}
        \end{column}
        \begin{column}{0.50\textwidth}
            \textbf{Métricas Cuantitativas de Impacto:}
            \centering
            \small
            \begin{tabular}{lcc}
                \toprule
                \textbf{Métrica de Rúbrica} & \textbf{Estático} & \textbf{Dinámico (IA)} \\
                \midrule
                Tiempo Medio Espera & 28.8 min & \textbf{12.4 min ($\downarrow$ 56.8\%)} \\
                Espera UniNorte P4 & 44.5 min & \textbf{14.8 min ($\downarrow$ 66.7\%)} \\
                Rechazo Sobrecupo & 39.4\% & \textbf{20.3\% ($\downarrow$ 48.6\%)} \\
                Rechazo UniNorte & 68.2\% & \textbf{15.5\% ($\downarrow$ 77.3\%)} \\
                \bottomrule
            \end{tabular}
            \vspace{0.2cm}
            \begin{exampleblock}{Comando de Ejecución Local}
                \texttt{python main.py} \quad $\implies$ \quad \texttt{http://localhost:8080}
            \end{exampleblock}
        \end{column}
    \end{columns}
\end{frame}

% -------------------------------------------------------------
% SLIDE 11: Conclusiones y Defensa ante Preguntas (Equipo)
% -------------------------------------------------------------
\begin{frame}{Conclusiones y Blindaje ante Preguntas de Cátedra}
    \textbf{Conclusiones Principales:}
    \begin{enumerate}
        \item Se modeló exitosamente la red vial en \textbf{doble escala}: subgrafo de 6 nodos con resultado conocido y red real de 15 paraderos con simulación de cierres por obras.
        \item El tratamiento axiomático de geodésicas empatadas ($\sigma_{st}(v)/\sigma_{st} = 0.5$) y la normalización combinatoria $\binom{5}{2}=10$ cuantifican rigurosamente el desvío.
        \item El Perceptrón de Rosenblatt resolvió la frontera en 2 épocas con error nulo, aportando rigor y transparencia matemática sin cajas negras.
        \item Impacto validado: reducción de $>56\%$ en tiempo de espera y $>48\%$ en sobrecupo, con dashboard interactivo para toma de decisiones en tiempo real.
    \end{enumerate}

    \vspace{0.2cm}
    \begin{center}
        \textbf{\large ¡Muchas gracias!} \\
        \textit{Sesión abierta a preguntas del profesor Iván Gómez.}
    \end{center}
\end{frame}

\end{document}
